\documentclass[10pt, aspectratio=169]{slidesiffmodelo}
\usepackage{metadados}

% ==========================================
% DADOS DO TRABALHO
% ==========================================
\title{Título do Trabalho: subtítulo do trabalho}

\author{
    \orcidlink{0009-0003-6724-4640}Pedro Henrique Rocha de Andrade\inst{1}, 
    % (ORCID é OPCIONAL) Informe mais autores:
    Fulano da Silva\inst{2},
}

\institute{
    \inst{1} Instituto Federal Fluminense \textit{Campus} Bom Jesus do Itabapoana\\
    % Para mais uma instituição, informe:
    \inst{2} Instituto Federal Fluminense \textit{Campus} Itaperuna\\
}

\begin{document}

% ==========================================
% Comece a escrever aqui! Não apague \titlepage
% ==========================================

% ==========================================
% CAPA
% ==========================================
\begin{frame}[t]
    \titlepage
\end{frame}

% ==========================================
% Sumário
% Item opcional, se quiser, remova todo o frame (\begin{frame}{Frame Title}\end{frame})
% ==========================================
\begin{frame}[t,small]
    \frametitle{Sumário}
    \small
    \tableofcontents[hideallsubsections,sectionstyle=show]
\end{frame}

% ==========================================
% Introdução 
% ==========================================
\section{Introdução}

% ==========================================
% Slide 1 da Introdução (Com imagem!)
% ==========================================
\subsection{Contexto}
\begin{frame}[t, allowframebreaks,small]
    \frametitle{Contexto}    
    \begin{textoimagem}
        \begin{itemize}
            \item Telescópios de grande porte;
            \item \textit{Big Data};
            \item Propriedades estelares;
            \item Características químicas e dinâmicas da vizinhança solar;
            \item Arqueologia Galáctica;
            \item Evolução da Via Láctea.
        \end{itemize}
        
    \colunaimagem
% ==========================================
% \inserirfigura[LARGURA (opcional, use para redimensionar)]{local do arquivvo}{Título}{Fonte (deixe vazio se prefrir.)}{Label para chamadas no texto} 
% ==========================================    
        \inserirfigura[width=0.8\linewidth]{img/gaia.jpg}{Satélite \emph{Gaia}, da Agência Espacial Europeia (ESA) Créditos: ESA/ATG medialab}{}{fig:gaia} 
    \end{textoimagem}
\end{frame}

% ==========================================
% Slide 2 da Introdução (Com imagem!)
% ==========================================
\subsection{Catálogos}
\begin{frame}[t, allowframebreaks]
    \frametitle{Surveys}
    \begin{textoimagem}
        \begin{itemize}
            \item Cruzamento de Dados;
            \item \emph{Gaia} Catalogue of Nearby Stars (GCNS), com $\sim 330$ mil \estrela até 100pc;
            \item Galactic Archaeology with HERMES (GALAH DR4), com $\sim 900$ mil \estrela.
        \end{itemize}
        
    \colunaimagem
% ==========================================
% \inserirfigura[LARGURA (opcional, use para redimensionar)]{local do arquivvo}{Título}{Fonte (deixe vazio se prefrir.)}{Label para chamadas no texto} 
% ==========================================
        \inserirfigura[width=0.8\linewidth]{img/aat_galah.jpeg}{Anglo-Australian Telescope. Créditos: Siding Spring Observatory.}{}{fig:galah} 
    \end{textoimagem}
\end{frame}

% ==========================================
% Objetivos 
% ==========================================
\section{Objetivos}
\begin{frame}[t]
    \frametitle{Objetivos}
    \begin{itemize}
        \item Caracterizar as propriedades físicas, químicas e cinemáticas das \estrela;
        \item Compreender a composição química;
        \item Evolução dinâmica da nossa vizinhança solar.
    \end{itemize}
\end{frame}

% ==========================================
% Metodologia 
% ==========================================
\section{Metodologia}
\begin{frame}[t, img]
\frametitle{Metodologia}
    \begin{itemize}
        \item Amostra final: $\sim 6\,000$ \estrela;
        \item Parâmetros da limpeza: 
        \begin{itemize}
            \item \feh;
            \item \mgfe; 
            \item \teff;
            \item \logg;
        \end{itemize}
    \end{itemize}
\end{frame}

% ==========================================
% Resultados 
% ==========================================
\section{Resultados}

% ==========================================
% Slide 1 de Resultado:
% ==========================================
\begin{frame}[t,img]
    \frametitle{Resultados}
    \begin{textoimagem}
        \begin{enumerate}
            \item Diagrama de Abundâncias Químicas;
            \item Diagrama de Kiel e Idade;
            \item Diagrama de Toomre;
            \item Diagrama Tinsley-Wallerstein.
        \end{enumerate}
        
    \colunaimagem
    
        \inserirfigura{img/anatomia.png}{Anatomia da Via Láctea. Créditos}{}{fig:anatomia} 
    \end{textoimagem}
\end{frame}

% ==========================================
% Slide 2 de Resultado:
% ==========================================
\subsection{Abundâncias Químicas}
\begin{frame}[t,img]
    \frametitle{Diagrama de Abundâncias Químicas}
    % Sem colunas, ajustamos apenas o width para dar um visual equilibrado no centro
    \vspace*{-1em}
    \inserirfigura[width=0.9\linewidth]{img/SBPC_Fig1_Feh_MgFe.pdf}{Distribuição das abundâncias de \feh \, e \mgfe.}{}{fig:histograma}
\end{frame}

% ==========================================
% Slide 3 de Resultado:
% ==========================================
\subsection{Kiel e Idade}
\begin{frame}[t,img]
    \frametitle{Diagrama de Kiel e Idade}
    \vspace*{-1em}
    \inserirfigura[width=0.9\linewidth]{img/SBPC_Fig2_Kiel_Idade.pdf}{Diagrama de Kiel e distribuição de idades estelares.}{}{fig:kiel}
\end{frame}

% ==========================================
% Slide 4 de Resultado:
% ==========================================
\subsection{Toomre}
\begin{frame}[t, img] % Ativa a fonte menor automaticamente
    \frametitle{Diagrama de Toomre}
    \begin{textoimagem}
        \begin{itemize}
            \item Círculos concêntricos em intervalos de 50 km s$^{-1}$;
            \item 228 \estrela prováveis pertences ao halo; \cite{Bovy2015}
            \item 5.970 \estrela prováveis pertencentes ao disco.
        \end{itemize}
    \colunaimagem
        \inserirfigura[width=0.9\linewidth]{img/SBPC_Fig3_Toomre.pdf}{Diagrama de Toomre.}{}{fig:toomre}
    \end{textoimagem}
\end{frame}

% ==========================================
% Slide 5 de Resultado:
% ==========================================
\subsection{Tinsley-Wallerstein}
\begin{frame}[t, img]
    \frametitle{Diagrama Tinsley-Wallerstein}
    \begin{textoimagem}
        \begin{itemize}
            \item \mgfe \, vs. \feh. Utilizado para separação do disco fino/espesso;
            \item \estrela abaixo da \textit{spline} são do disco fino, acima, do espesso;
            \item A opacidade reflete a incerteza na classificação. 
        \end{itemize}
        
    \colunaimagem
    \inserirfigura[width=0.9\linewidth]{img/SBPC_Fig4_Tinsley_Wallerstein.pdf}{Diagrama de Tinsley-Wallerstein.}{}{fig:tinsley}
    \end{textoimagem}
\end{frame}

% ==========================================
% Conclusões e Perspectivas
% ==========================================
\section{Conclusões e Perspectivas}
\begin{frame}[t, allowframebreaks]
    \frametitle{Conclusões e Perspectivas}
    \begin{itemize}
        \item \estrela predominantemente do tipo espectral parecidas com o Sol;
        \item Idade mediana $\approx 1,6$ bilhões de anos (Ga);
        \item Ligeira deficiência metálica ($\feh \approx -0,19$ dex.);
        \item Concordância de espectroscopia e isócronas teóricas;
        \item Dominância de \estrela do disco, poucas \estrela no halo;
        \item As \estrela não são tão parecidas com o Sol quanto se esperava!
    \end{itemize}
\end{frame}

% ==========================================
% Slide Final
% ==========================================
\begin{frame}[t]
    \begin{center}
        \vfill
        {\Huge \textbf{Obrigado!}} \\
        \vspace*{1em}
        % insira aqui seu email ou contato 
    \end{center}
\end{frame}

% ==========================================
% Agradecimentos (Opcional, caso queira, apague todo o frame)
% ==========================================
\section*{Agradecimentos}
\begin{frame}[t,small]
    \frametitle{Agradecimentos}
    \small
    \begin{itemize}
        \item Agradecemos...
    \end{itemize}   
\end{frame}

% ==========================================
% Referências Bibliográficas
% Indique o arquivo referencias.bib, e cole todas as suas citações em formato bibtex lá. As citações são automáticas quando você as chama no texto por \citeonline{} ou \cite{}. Se precisa forçar aparecer, use \nocite{} aqui.
% ==========================================
\section*{Referências}
\begin{frame}[allowframebreaks]
    \frametitle{Referências}   
    \bibliography{referencias}
\end{frame}

\end{document}